%% ============================================================
%%  CHAPTER 1 — INTRODUCTION
%% ============================================================
\chapter{Introduction}
\label{ch:introduction}

\section{Background}
\label{sec:background}

Upper-limb amputation represents one of the most debilitating physical
conditions a person can endure, fundamentally disrupting independence
and quality of life across virtually every daily activity — from eating
and writing to personal hygiene and social interaction. Globally,
millions of individuals live with upper-limb loss resulting from
traumatic injuries, congenital conditions, vascular disease, and
conflict-related incidents~\cite{ziegler2008prevalence}. In
conflict-affected regions such as Palestine, decades of violence have
produced a significant population of young individuals with traumatic
upper-limb amputations, often accompanied by extensive peripheral nerve
and soft tissue damage that severely limits the effectiveness of
conventional prosthetic solutions~\cite{who2022assistive}.

Modern prosthetic hands have advanced considerably, incorporating
multiple degrees of freedom, anthropomorphic designs, and sophisticated
actuation mechanisms~\cite{ten2017review}. However, the utility of
any prosthetic device is ultimately constrained by the quality of its
control interface. The dominant technology for non-invasive prosthetic
control is myoelectric control, which decodes surface
electromyography (\acrshort{emg}) signals from residual muscles in the
residual limb. While effective for many amputees, myoelectric control
fails entirely when the residual musculature is insufficient, when
peripheral nerve damage prevents reliable muscle activation, or when
the amputation is at a high level — conditions that are precisely
characteristic of war-related injuries~\cite{al2013classification}.

Brain--Computer Interface (\acrshort{bci}) technology addresses this
fundamental limitation by enabling direct communication between the
brain and an external device, completely bypassing the damaged
peripheral neuromuscular pathway~\cite{wolpaw2002bci}. Rather than
relying on residual muscle activity, a \acrshort{bci} reads electrical
signals directly from the scalp and translates them into control
commands through signal processing and machine learning algorithms.
Among the various non-invasive \acrshort{bci} modalities,
Electroencephalography (\acrshort{eeg}) is the most widely adopted,
offering a practical balance between signal quality, cost,
portability, and safety~\cite{lotte2007review}.

Early \acrshort{eeg}-based \acrshort{bci} systems for prosthetic
control focused predominantly on \textit{motor imagery} --- the
mental simulation of physical movements. When a person imagines
moving their hand, characteristic modulations in the mu
(\SIrange{8}{13}{\hertz}) and beta (\SIrange{13}{30}{\hertz})
sensorimotor rhythms over the motor cortex can be detected and
classified~\cite{pfurtscheller1999erd}. However, motor imagery
\acrshort{bci} systems face well-documented practical
challenges: the neural signals are subtle, highly variable across
individuals and sessions, require extensive subject-specific
calibration, and demand sophisticated spatial filtering techniques
such as Common Spatial Patterns (\acrshort{csp}) to extract
meaningful features~\cite{blankertz2008csp}. Classification accuracies
in real-world deployments rarely match laboratory benchmarks, and a
significant proportion of users — sometimes referred to as
\textit{BCI illiterates} — cannot reliably modulate their sensorimotor
rhythms at all~\cite{vidaurre2010bci}.

An alternative and increasingly attractive paradigm exploits
\textit{voluntary physiological artifacts} — large-amplitude
\acrshort{eeg} signals produced by intentional facial and jaw muscle
activations, such as eye blinks, jaw clinching, and bruxism (lateral
teeth grinding). Although traditionally treated as nuisances to be
removed during preprocessing, these artifact signals possess
characteristics that are arguably more suitable for practical
\acrshort{bci} control than motor imagery: they are orders of
magnitude larger in amplitude (tens to hundreds of microvolts versus
single-digit microvolts for motor imagery), highly reproducible within
a subject, distinguishable by their spatial and temporal signatures,
and can be performed discretely and deliberately without attracting
social attention~\cite{usakli2010artifact}. The key insight is that
these signals, while artifactual from a neuroscience perspective, are
entirely \textit{voluntary} from a control perspective — the user
intentionally produces them to issue commands.

It is worth noting at the outset that these voluntary artifacts are
not, strictly speaking, cortical signals. The blink commands are
ocular potentials (electrooculography, \acrshort{eog}) generated by
eye movement and captured at the frontal electrodes, while the jaw
commands are muscle potentials (\acrshort{emg}) generated by the
muscles of mastication and captured at the temporal electrodes. They
are recorded by a scalp \acrshort{eeg} montage and processed by a
\acrshort{bci} pipeline, but their physiological origin is ocular and
muscular rather than purely cortical. This dual signal nature, far
from being a weakness, is exploited deliberately in the final system
design (Chapter~9), where each signal family is routed to the
classification method best matched to its physiology. The
neurophysiological basis of these signals is developed in detail in
Chapter~3.

\section{Problem Definition}
\label{sec:problem}

This project addresses the following core problem:

\begin{quote}
\textit{How can a reliable, real-time, non-invasive \acrshort{bci}
system be designed and implemented to provide five distinct prosthetic
hand control commands using voluntary \acrshort{eeg} artifact
signals, with classification accuracy and latency suitable for
practical assistive device deployment?}
\end{quote}

Several specific technical challenges must be overcome:

\begin{enumerate}[leftmargin=*, label=\textbf{P\arabic*.}]

    \item \textbf{Heterogeneous signal types.} The five commands are
    not produced by a single physiological mechanism. Blink commands
    are ocular (\acrshort{eog}) signals at the frontal electrodes,
    whereas clinch and bruxism commands are muscular (\acrshort{emg})
    signals at the temporal electrodes. A feature that is highly
    discriminative for one family (e.g. muscle energy) is pure noise
    for the other, so a single monolithic classifier cannot apply
    features conditionally — a fundamentally architectural challenge.

    \item \textbf{Blink sub-classification under label contamination.}
    Distinguishing single from double blinks is intrinsically hard:
    a substantial fraction of recorded ``single blink'' trials in
    practice contain rapid clusters of multiple blinks, contaminating
    the class labels. Event-counting methods fail under this
    contamination, demanding a shape-based rather than count-based
    approach.

    \item \textbf{Cross-session amplitude drift.} \acrshort{eeg}
    and \acrshort{emg} signals exhibit non-stationarity — their
    absolute amplitude shifts between recording sessions due to
    electrode placement, impedance changes, and physiological state.
    Amplitude-based features that work within a session fail across
    sessions, requiring representations that are invariant to
    amplitude scaling.

    \item \textbf{Resting window contamination.} In any recording,
    active gesture epochs are interspersed with rest periods during
    which no command is intended. These rest windows carry the label
    of the surrounding gesture class, creating conflicting training
    examples that degrade performance if not removed.

    \item \textbf{Data leakage in evaluation.} When overlapping
    epochs are split randomly into training and test sets, temporally
    adjacent near-duplicate epochs appear in both sets, producing
    artificially inflated accuracy estimates. Honest evaluation
    requires a principled, recording-grouped, cross-session split
    strategy that prevents this leakage.

    \item \textbf{Real-time inference with safety guarantees.}
    Translating classifier outputs into prosthetic commands in
    real time requires not only low latency but also robust safety
    mechanisms that prevent unintended activations, filter transient
    misclassifications, and ensure deterministic, predictable hand
    behaviour.

\end{enumerate}

\section{Motivation}
\label{sec:motivation}

The motivation for this project arises from two intersecting
dimensions: technical limitations in existing \acrshort{bci} research
and pressing humanitarian needs.

\subsection{Technical Motivation}

Despite two decades of research, \acrshort{eeg}-based \acrshort{bci}
prosthetic control remains largely confined to laboratory settings.
The gap between offline classification accuracy and real-world
system reliability is well documented~\cite{vidaurre2010bci}, and
the dominant motor imagery paradigm faces inherent limitations in
signal amplitude, user-to-user variability, and training burden.
Artifact-based \acrshort{bci} represents an underexplored but
practically promising alternative, particularly for applications
requiring a small number of discrete, reliable commands — precisely
the requirement for prosthetic hand control. By leveraging signals
that are naturally large, spatially distinct, and easily reproducible,
an artifact-based system can potentially achieve higher real-world
reliability with simpler hardware and shorter calibration sessions.

\subsection{Clinical and Humanitarian Motivation}

Palestine and similar conflict-affected regions face a disproportionate
burden of traumatic upper-limb amputations. War-related amputations
frequently involve not only limb loss but also extensive tissue trauma,
nerve damage, and compromised residual musculature, rendering
myoelectric prostheses ineffective for a substantial portion of
affected individuals~\cite{who2022assistive}. Advanced commercial
prosthetic systems cost tens of thousands of dollars and require
specialised clinical infrastructure for fitting and maintenance —
resources that are scarce or unavailable in conflict zones. A
\acrshort{bci}-controlled prosthesis built from accessible, low-cost
components (consumer-grade \acrshort{eeg} hardware, open-source
software, 3D-printed mechanical parts) represents a meaningful step
toward democratising assistive technology for this population.

\section{Objectives}
\label{sec:objectives}

The primary objectives of this project are as follows:

\begin{enumerate}[leftmargin=*, label=\textbf{O\arabic*.}]

    \item \textbf{Design and implement a five-command artifact-based
    \acrshort{bci} pipeline} that classifies voluntary
    \acrshort{eeg} artifact signals — single and double eye blinks,
    jaw clinch, left bruxism, and right bruxism — into discrete
    prosthetic hand control commands in real time.

    \item \textbf{Develop a robust signal processing pipeline}
    encompassing \acrshort{edf} file loading, notch and bandpass
    filtering, Common Average Reference (\acrshort{car}) spatial
    normalisation, and activity-gated epoch extraction that
    suppresses resting-window contamination.

    \item \textbf{Engineer a physiologically motivated, channel-aware
    feature extraction scheme} producing a channel-aware feature
    superset from frontal and temporal electrode groups, incorporating
    time-domain, frequency-domain, and asymmetry features tailored
    to the specific artifact types being classified.

    \item \textbf{Design a hybrid, physiology-matched classification
    architecture} that routes each gesture family to the method best
    suited to its signal type — hand-crafted features for the
    blink-versus-muscle separation, Dynamic Time Warping
    (\acrshort{dtw}) shape-matching for blink sub-classification, and
    Riemannian covariance geometry for muscle sub-classification —
    and validate it under a principled, leakage-free,
    Leave-One-Session-Out (\acrshort{loso}) evaluation protocol.

    \item \textbf{Implement a real-time inference engine} with a
    sliding window, activity gate, confidence threshold, and Finite
    State Machine (\acrshort{fsm}) that translates classifier outputs
    into safe, validated commands transmitted to an Arduino-controlled
    3D-printed prosthetic hand.

    \item \textbf{Design, fabricate, and integrate a 3D-printed
    prosthetic hand} capable of executing the five commanded actions
    (open, close, rotate left, rotate right, and confirm) driven by
    servo motors under Arduino firmware control.

    \item \textbf{Collect a dedicated, multi-session dataset} of
    voluntary \acrshort{eeg} artifact recordings in \acrshort{edf}
    format, covering all gesture classes across three independent
    recording sessions, to serve as the training and evaluation
    corpus.

\end{enumerate}

\section{Project Scope}
\label{sec:scope}

This project focuses on single-subject, subject-specific
\acrshort{bci} design — a deliberate choice consistent with the
intended deployment scenario, where a prosthetic device is
personalised to a single user. The system is not designed for
cross-subject generalisation; instead, it prioritises within-subject
reliability and session-to-session robustness for the same individual.

The control paradigm is \textit{discrete command-based}: each
voluntary artifact gesture maps to one of five fixed commands, and
the prosthetic hand transitions between a finite set of states
accordingly. Continuous proportional control — where the degree of
hand opening or wrist rotation is modulated by signal amplitude —
is explicitly outside the scope of this work, as it exceeds the
reliable information transfer capacity of artifact-based
\acrshort{eeg} signals.

The project does not include clinical trials with amputee subjects,
formal usability studies, regulatory approval processes, or
long-term home deployment. It demonstrates technical feasibility
and engineering proof-of-concept, establishing a foundation upon
which future clinical research may build.

\section{Key Contributions}
\label{sec:contributions}

This project makes the following original contributions:

\begin{enumerate}[leftmargin=*, label=\textbf{C\arabic*.}]

    \item \textbf{Physiology-matched hybrid classification
    architecture.} The central contribution is a three-way hybrid
    classifier that routes each gesture family to the method best
    matched to its underlying signal physiology: hand-crafted
    channel-aware features perform the blink-versus-muscle separation,
    Dynamic Time Warping (\acrshort{dtw}) shape-matching
    sub-classifies the blink commands, and Riemannian covariance
    geometry sub-classifies the muscle commands. This architecture
    achieves 91.3\% accuracy under the strict cross-session
    \acrshort{loso} protocol, compared with 85.1\% for a flat
    Random Forest baseline.

    \item \textbf{Riemannian geometry for amplitude-drift-robust
    muscle classification.} The use of inter-channel covariance
    matrices in a Riemannian tangent space, rather than absolute
    amplitude features, is shown to survive the day-to-day amplitude
    drift that defeats amplitude-based descriptors — improving overall
    accuracy by +4.1\% and lifting the muscle-class F1 scores above
    0.95.

    \item \textbf{DTW shape-matching for contamination-robust blink
    classification.} Where five separate event-counting methods
    failed on contaminated single-blink data, \acrshort{dtw}
    shape-matching with $k$-nearest-neighbour exemplar voting succeeds
    by comparing full waveform morphology rather than counting events
    — improving overall accuracy by a further +2.1\% and lifting the
    blink-class F1 scores.

    \item \textbf{Activity-gated epoch extraction.} A preprocessing
    step that uses per-recording \acrshort{rms} energy thresholding
    on physiologically relevant electrode pairs (Fp1/Fp2 for blinks,
    T3/T4 for jaw) to automatically reject resting windows mislabelled
    with gesture class labels, yielding a +3.1\% improvement in
    classification accuracy.

    \item \textbf{Rigorous, leakage-free evaluation methodology.}
    A systematic data-leakage audit — including a label-permutation
    shuffle test that confirms the pipeline does not manufacture
    signal — demonstrates that naive random splitting of overlapping
    epochs inflates accuracy by 8--10 points. All headline results
    are reported under the strict, deployment-realistic
    Leave-One-Session-Out (\acrshort{loso}) protocol.

    \item \textbf{Real-time inference engine with FSM safety layer.}
    A production-ready inference pipeline that mirrors training
    preprocessing exactly, incorporates an activity gate and
    confidence threshold to suppress false activations, and uses a
    Finite State Machine with a three-window confirmation requirement
    to ensure safe, deterministic prosthetic behaviour at a per-window
    latency of \SI{14}{\milli\second}.

    \item \textbf{Self-collected multi-session dataset and low-cost
    integrated system.} A purpose-built dataset of voluntary artifact
    recordings across three independent sessions using a 19-channel
    \acrshort{eeg} device at \SI{200}{\hertz}, together with a complete
    end-to-end system from scalp signal to physical prosthetic
    actuation, implemented entirely with open-source software and a
    3D-printed prosthetic hand — demonstrating practical accessibility
    for resource-constrained deployment contexts.

\end{enumerate}

\section{Report Organisation}
\label{sec:organisation}

The remainder of this report is organised as follows:

\textbf{Chapter~2} reviews existing literature on
\acrshort{eeg}-based \acrshort{bci} systems for prosthetic control,
artifact-based \acrshort{bci} paradigms, signal processing and
classification methods, safety mechanisms, and 3D-printed prosthetic
technologies, identifying the research gaps that this project addresses.

\textbf{Chapter~3} establishes the medical and physiological
foundations of the artifact-based paradigm, detailing the
neurophysiology of the eye-blink and jaw artifacts, the anatomical
basis of left/right bruxism asymmetry, and why the relevant cranial
pathways remain intact in the target amputee population.

\textbf{Chapter~4} describes the system hardware, including the
\acrshort{eeg} acquisition device and electrode configuration, the
3D-printed prosthetic hand design and fabrication, and the
Arduino-based control interface.

\textbf{Chapter~5} provides a high-level methodology overview,
walking through the full processing pipeline stage by stage with
diagrams before the subsequent chapters examine each stage in detail.

\textbf{Chapter~6} details the experimental protocol and data
collection methodology, covering the recording setup, session
structure, artifact gesture protocols, dataset statistics, and the
data quality issues discovered.

\textbf{Chapter~7} presents the signal processing pipeline, covering
each preprocessing stage from \acrshort{edf} loading through notch
and bandpass filtering, \acrshort{car} normalisation, and
activity-gated epoch extraction, with full mathematical derivations.

\textbf{Chapter~8} describes the channel-aware feature extraction
scheme, detailing all feature categories, their physiological
justifications, the T3--T4 asymmetry index, and the systematic
feature engineering experiments that motivated the architectural
approach.

\textbf{Chapter~9} covers the classification methodology in full:
the evaluation-methodology and data-leakage audit, the label
taxonomy evolution, and the three-way hybrid architecture combining
hierarchical routing, Riemannian muscle classification, and
\acrshort{dtw} blink classification.

\textbf{Chapter~10} presents the experimental results and analysis,
including the performance progression, per-class metrics, confusion
matrix analysis, and an ablation study quantifying the contribution
of each design decision.

\textbf{Chapter~11} describes the configuration and control
application --- the offline, controller-hosted web app through which
the user remaps artifact commands to hand actions, creates custom
poses, and chains them into sequences.

\textbf{Chapter~12} describes the real-time inference system,
detailing the sliding window mechanism, activity gate, confidence
thresholding, \acrshort{fsm} design, the UART command protocol, and
end-to-end latency analysis.

\textbf{Chapter~13} discusses the broader implications of the
results, the advantages and limitations of the artifact-based
paradigm, and the ethical considerations of deploying \acrshort{bci}
technology in assistive contexts.

\textbf{Chapter~14} concludes the report with a summary of
contributions, lessons learned, and directions for future work.